\documentclass[11pt]{article}
\usepackage{varnernotes}

\begin{document}

\lecturenotes{15b}{Launching the 2026 Autotrader}{Fall 2026}{December 3, 2026}{Jeffrey D. Varner}

\learningobjectives
  {Specify the decision pipeline}
  {Connect lagged data, EMA sentiment, single-index forecasts, positive utility preferences, target holdings, orders, fills, and ledger updates.}
  {Enforce launch invariants}
  {Test information timing, numerical domains, self-financing, inventory, cash, execution, and risk limits before any external order is possible.}
  {Stage an evidence-based launch}
  {Progress from deterministic simulation through walk-forward and paper trading to a capped pilot with monitoring, reconciliation, and a kill switch.}

\section{A Trading System Is More Than a Signal}

An autotrader is a stateful control system coupled to market data, a broker, and an accounting ledger. The model proposes an action; the execution system determines what was actually filled; the ledger and risk engine determine what may happen next. A backtest that jumps directly from a signal to portfolio return omits the interfaces where many costly failures occur.

For the course prototype, a market signal comes from short- and long-horizon exponential moving averages (EMAs). With price observation $P_t$ and span $N$, define
\begin{equation}
  \label{eq:ema}
  E_t^{(N)}=\alpha_NP_t+(1-\alpha_N)E_{t-1}^{(N)},
  \qquad
  \alpha_N:=\frac{2}{N+1}.
\end{equation}
The initialization rule for $E_0^{(N)}$ must be declared because it affects early signals. A dimensionless crossover score available after the close of $t$ is
\begin{equation}
  \label{eq:sentiment}
  z_t:=\frac{E_t^{(s)}}{E_t^{(\ell)}}-1,
  \qquad s<\ell,
\end{equation}
where $s$ and $\ell$ are the short and long spans. The notebook uses $-Gz_t$ to tilt allocation, where $G>0$ is a gain. Because this quantity may be negative, it is better called a \emph{sentiment or exposure tilt} than classical risk aversion, which is normally nonnegative.

\begin{figure}[ht]
  \centering
  \begin{tikzpicture}[>=Latex,font=\sffamily\small,node distance=0.78cm]
    \node[draw=vnink,very thick,rounded corners,fill=vnsoft,minimum width=1.62cm,minimum height=0.85cm,align=center] (data) {validated\\data};
    \node[draw=vnink,very thick,rounded corners,fill=white,minimum width=1.62cm,minimum height=0.85cm,align=center,right=of data] (signal) {lagged signal\\and forecast};
    \node[draw=vnink,very thick,rounded corners,fill=vnsoft,minimum width=1.62cm,minimum height=0.85cm,align=center,right=of signal] (target) {risk-checked\\target};
    \node[draw=vnink,very thick,rounded corners,fill=white,minimum width=1.62cm,minimum height=0.85cm,align=center,right=of target] (order) {orders and\\fills};
    \node[draw=vnink,very thick,rounded corners,fill=vnsoft,minimum width=1.62cm,minimum height=0.85cm,align=center,right=of order] (ledger) {reconciled\\ledger};
    \draw[vnlink,very thick,->] (data)--(signal);
    \draw[vnlink,very thick,->] (signal)--(target);
    \draw[vnlink,very thick,->] (target)--(order);
    \draw[vnlink,very thick,->] (order)--(ledger);
    \draw[vncarnelian,very thick,->,bend left=22] (ledger) to node[below,align=center] {positions, cash, and risk\\become next state} (data);
  \end{tikzpicture}
  \caption{The minimum auditable autotrader loop. Each arrow is an interface with timing, validation, failure, and recovery requirements.}
  \label{fig:autotrader-loop}
\end{figure}

If the signal uses the close at date $t$, the first routine execution is at a price after that close, such as the next session's open. Same-close execution is valid only when the signal was computable before the order deadline from information already observed. Timestamp every observation, feature, decision, order, acknowledgement, fill, and ledger event.

\section{Forecasts and Preferences Must Stay in Domain}

For ticker $i$, a single-index forecast can be written
\begin{equation}
  \label{eq:sim-forecast}
  \widehat g_{i,t}=\widehat\alpha_{i,t}
  +\widehat\beta_{i,t}\widehat g_{m,t},
\end{equation}
where the estimates and market forecast use data available at the decision time. A separate transformation converts this forecast and the sentiment score into a strictly positive allocation preference $\gamma_{i,t}>0$. One stable example is
\begin{equation}
  \label{eq:positive-preference}
  \gamma_{i,t}:=\gamma_{\min}+
  \log\!\left(1+\exp\!\left[
    a\widehat g_{i,t}-b\rho_{i,t}+cz_t
  \right]\right),
\end{equation}
where $\gamma_{\min}>0$ is a preference floor, $a,b,c$ are declared scaling coefficients, and $\rho_{i,t}\geq0$ is a positive risk measure such as forecast volatility. Standardize the inputs or bound the score to avoid overflow and unstable concentration.

This construction separates expected reward from risk. A SIM beta $\widehat\beta_i$ is signed market exposure, not a universally positive risk scale. The inherited prototype computes $\widehat\beta_i^{\lambda_t}$ for a generally non-integer, signed tilt $\lambda_t$; this is not real-valued when $\widehat\beta_i<0$ and is undefined at zero for some exponents. Use a positive $\rho_i$, or explicitly transform magnitude and sign separately.

Likewise, $\tanh(\widehat g_i)$ lies in $(-1,1)$ and can produce negative values. Those values cannot be inserted as exponents into the positive Cobb--Douglas solution derived in L14a. Equation \eqnref{positive-preference}, explicit exclusion of non-preferred assets, or a different constrained objective fixes the domain mismatch. A tiny forced holding of every negative-score asset is a policy choice, not the analytical optimum.

With positive preferences, normalized target weights are
\begin{equation}
  \label{eq:target-weights}
  w_{i,t}^*=\frac{\gamma_{i,t}}{\sum_{j\in\mathcal S_t}\gamma_{j,t}},
  \qquad
  n_{i,t}^*=\frac{w_{i,t}^*B_t}{p_{i,t}^{\mathrm{ref}}},
\end{equation}
before implementability constraints. Here $\mathcal S_t$ is the point-in-time eligible universe, $B_t$ is deployable capital, and $p_{i,t}^{\mathrm{ref}}>0$ is a reference price used only to form targets. The execution engine must recalculate feasibility using actual order sides and fills.

\begin{remark}{Synthetic realism is a test claim, not an input label}{synthetic-realism}
An HMM-driven market factor combined with SIM residuals can generate useful controlled scenarios, but calibration to historical data does not make simulated paths realistic by declaration. Check regime duration, cross-sectional covariance, tails, volatility clustering, jumps, and drawdowns, and stress features the generator does not reproduce.
\end{remark}

\notebooklink
  {L15b: Synthetic Trade Bot}
  {https://github.com/varnerlab/CHEME-5660-CourseRepository-Fall-2026/blob/main/lectures/week-15/L15b/CHEME-5660-L15b-Example-MyTradeBot-Synthetic-Fall-2025.ipynb}
  {Generate synthetic HMM/SIM prices, form EMA crossover signals, allocate shares, and compare the prototype with SPY and STRIPS benchmarks.}

\section{The Ledger Is the Source of Truth}

Let $h_{i,t}^-$ and $c_t^-$ be broker-confirmed shares and cash before an execution cycle. If signed fill $f_{i,t}$ is positive for a purchase and negative for a sale, post-fill state is
\begin{align}
  \label{eq:fill-ledger}
  h_{i,t}^+&=h_{i,t}^-+f_{i,t},\\
  c_t^+&=c_t^--\sum_i p_{i,t}^{\mathrm{fill}}f_{i,t}
  -C_t(\mathbf f_t),
\end{align}
where $p_{i,t}^{\mathrm{fill}}>0$ and total execution cost $C_t\geq0$ includes commissions, fees, spread, and modeled impact. Marked wealth at a declared liquidation convention $p_{i,t}^{\mathrm{mark}}$ is
\begin{equation}
  \label{eq:marked-wealth}
  W_t=c_t^++\sum_i h_{i,t}^+p_{i,t}^{\mathrm{mark}}.
\end{equation}
Order submission does not change the ledger; fills do. Partial fills, rejects, cancellations, and duplicate messages must be handled explicitly.

\begin{theorem}{Execution conservation check}{execution-invariant}
Ignoring external deposits, withdrawals, dividends, and mark changes during the fill event, \eqnref{fill-ledger} implies
\begin{equation}
  \label{eq:execution-conservation}
  \left(c_t^++\sum_i h_{i,t}^+p_{i,t}^{\mathrm{fill}}\right)
  -\left(c_t^-+\sum_i h_{i,t}^-p_{i,t}^{\mathrm{fill}}\right)
  =-C_t(\mathbf f_t).
\end{equation}
Any larger discrepancy signals a sign, multiplier, quantity, price, fee, or event-processing error and should stop the strategy before another order cycle.
\end{theorem}

The notebook's pedagogical rule liquidates at a daily price and buys at a random perturbation of another price. For research, fill simulation should be side-aware and feasible: buys at or above the ask, sells at or below the bid, participation limited by volume, and fills bounded to positive prices. A Gaussian percentage perturbation has unbounded support and does not by itself model spreads or market impact.

Every cycle should assert invariants before and after trading:
\begin{itemize}[leftmargin=1.3em,itemsep=2pt,topsep=2pt]
  \item all prices, quantities, timestamps, and identifiers are valid and finite;
  \item the universe, corporate actions, and market session are current; stale or missing data cause no trade;
  \item cash, shorting, leverage, per-name weight, gross/net exposure, and turnover obey configured limits;
  \item order IDs are idempotent, open orders are included in risk, and broker state reconciles to the internal ledger; and
  \item a breached drawdown, connectivity, data-quality, or reconciliation threshold cancels open orders and disables new ones.
\end{itemize}

\section{Launch Through Evidence Gates}

The 2026 Autotrader should move through progressively more realistic environments. Advancement is conditional on predeclared evidence, not on one attractive wealth chart.

\begin{figure}[ht]
  \centering
  \begin{tikzpicture}[>=Latex,font=\sffamily\small,node distance=0.7cm]
    \foreach \name/\text/\x in {u/{unit and\\property tests}/0,s/{deterministic\\simulation}/2.35,w/{walk-forward\\backtest}/4.7,p/{paper and\\shadow}/7.05,l/{capped\\pilot}/9.4}{
      \node[draw=vnink,very thick,rounded corners,fill=vnsoft,minimum width=1.95cm,minimum height=0.92cm,align=center] (\name) at (\x,0) {\text};
    }
    \draw[vnlink,very thick,->] (u)--(s);
    \draw[vnlink,very thick,->] (s)--(w);
    \draw[vnlink,very thick,->] (w)--(p);
    \draw[vncarnelian,very thick,->] (p)--(l);
    \node[below=0.6cm of w,align=center,text=vnmuted] {reconcile, review, and require a pass at every gate};
  \end{tikzpicture}
  \caption{A safe launch ladder. Live capital appears only in the final capped pilot, after the identical decision and ledger code has passed offline and broker-connected tests.}
  \label{fig:launch-ladder}
\end{figure}

\textbf{Unit and property tests} cover EMA initialization, strict lags, preference-domain checks, allocation constraints, rounding, fill signs, conservation \eqnref{execution-conservation}, duplicate events, and deterministic seeds. \textbf{Simulation} adds regime, gap, spread, missing-data, rejection, partial-fill, and disconnect scenarios. \textbf{Walk-forward testing} repeatedly freezes parameters, trades a later interval, and reports every trial rather than selecting the best seed or window.

\textbf{Paper trading} exercises real market-data and broker interfaces without economic fills. A parallel shadow ledger compares intended orders, broker acknowledgements, hypothetical fills, and end-of-day positions. Only after sustained reconciliation and operational review should a \textbf{capped pilot} permit a small approved universe, capital limit, maximum order size, rate limit, loss limit, and staffed kill switch. Whether any live pilot is appropriate is a separate governance decision, not a course deliverable.

Performance reporting uses the ledger and includes return, volatility, maximum drawdown, expected shortfall, turnover, costs, exposure, rejected and partial orders, stale-data events, reconciliation breaks, and downtime. Compare with passive SPY, the declared cash/STRIPS alternative, equal weight, and no-rebalance variants under the same dates and capital convention. Synthetic success, historical backtest success, and paper success are distinct evidence tiers.

\keytakeaways
  {In this lecture, we converted the course prototype into an auditable launch design with valid preference mathematics, execution accounting, risk controls, and evidence gates.}
  {Stay inside the model domain}
  {Signed SIM beta is not a positive risk scale, signed sentiment is not classical risk aversion, and Cobb--Douglas preferences must remain positive.}
  {Fills change the account}
  {The broker-confirmed ledger, conservation checks, open-order exposure, and reconciliation---not target weights or submitted orders---define holdings, cash, and wealth.}
  {Launch is progressive}
  {Deterministic tests, stress simulation, walk-forward evaluation, paper/shadow operation, and governance precede any tightly capped live pilot.}
  {The course ends with a system that can be questioned, reproduced, monitored, and improved---the essential traits of responsible quantitative engineering.}

\begin{readinglist}
  \item Marcos L\'opez de Prado, \emph{Advances in Financial Machine Learning}, Wiley (2018), chapters on backtest overfitting and feature importance.
  \item David H. Bailey et al., \href{https://doi.org/10.1093/rfs/hhu041}{``The Probability of Backtest Overfitting,''} \emph{Journal of Computational Finance} 20(4), 39--69 (2017).
  \item Philip Maymin, \href{https://doi.org/10.3905/jot.2011.6.1.001}{``Markets Are Not Continuous,''} \emph{Journal of Trading} 6(1), 1--8 (2011).
\end{readinglist}

{\sffamily\footnotesize\color{vnmuted}\textbf{Educational use and risk notice.} These notes are for instruction only. Automated trading can lose money rapidly and can amplify software, data, execution, and model errors.}

\end{document}
